% !TeX root = ../../main.tex
% =====================================================================
%  Casos de carga y dimensionado
%  Responsable:
%  Última revisión:
% =====================================================================
\chapter{Casos de carga y dimensionado}
\label{ch:casos_carga}

\begin{objetivos}
  \item Definir una tabla de casos de carga defendible, sabiendo de dónde sale
        cada número.
  \item Justificar los factores de bache y los coeficientes de seguridad en vez
        de heredarlos.
  \item Resolver a mano el equilibrio de una esquina y obtener la carga de cada
        barra.
  \item Usar el elemento finito para lo que sirve, con convergencia de malla
        demostrada.
  \item Decidir qué merece la pena ensayar físicamente y con qué montaje.
\end{objetivos}

\fichasesion{120 minutos}{Capítulos \ref{ch:mecanica_aplicada},
\ref{ch:regimen_estacionario} y \ref{ch:cinematica_suspension}}%
{Ordenador con \texttt{Codigo/suspension/cargas\_esquina.py} y el CAD}

\begin{clave}
La pregunta que más suspende en el design event es \textbf{«¿de dónde salen tus
cargas?»}. No se contesta con un coeficiente de seguridad ni con un color de
tensiones: se contesta con una cadena que empieza en una aceleración medida o
razonada, pasa por la transferencia de carga del capítulo
\ref{ch:regimen_estacionario} y termina en la fuerza axial de cada barra. Este
capítulo es esa cadena.
\end{clave}

% ---------------------------------------------------------------------
\section{Definición de los casos de carga}
\label{sec:casos_carga:01}
\index{casos de carga}

\subsection{De dónde salen las aceleraciones}

Un caso de carga es una combinación de aceleraciones que el coche va a ver.
Hay cuatro sitios de donde pueden salir, y no valen lo mismo:

\begin{center}
\begin{tabularx}{\textwidth}{@{}lXl@{}}
\toprule
\textbf{Origen} & \textbf{Qué tan bueno es} & \textbf{Cuándo} \\
\midrule
Telemetría propia & Lo mejor: es tu coche en tu pista & Cuando ya se ha rodado \\
Modelo de neumático propio & Muy bueno si el modelo está validado & Cap. \ref{ch:modelado_neumatico} \\
Datos de equipos parecidos & Aceptable como punto de partida & Primer coche \\
Valores de la bibliografía & El último recurso, y hay que decirlo & Primer coche \\
\bottomrule
\end{tabularx}
\end{center}

Los valores de referencia que se manejan en Formula Student ---y que hay que
usar como punto de partida, no como respuesta--- están en torno a
\SIrange{1.5}{2.5}{\gacc} laterales según neumático y aerodinámica,
\SIrange{1.5}{2}{\gacc} de frenada y \SIrange{3}{5}{\gacc} verticales en bache
o bordillo.

\begin{ojo}
\textbf{Los nuestros tienen que salir de nuestros datos, no de esa lista.} Un
coche sin aerodinámica y con un neumático duro no da \SI{2.5}{\gacc}: dar por
buena la cifra alta de la horquilla produce una suspensión pesada dimensionada
para un coche que no existe. Y al revés, usar la cifra baja sin justificarla es
igual de indefendible.
\end{ojo}

\subsection{La tabla}

Los casos que hay que tener, como mínimo, son estos. La columna vertical incluye
el factor de bache del apartado siguiente:

\begin{center}
\begin{tabularx}{\textwidth}{@{}lSSSX@{}}
\toprule
\textbf{Caso} & {$\ax$ (\si{\gacc})} & {$\ay$ (\si{\gacc})} & {$a_z$ (\si{\gacc})} & \textbf{Qué dimensiona} \\
\midrule
Frenada máxima      & -1.5 & 0.0 & 1.0 & Trapecios, anclajes, frenos \\
Curva máxima        &  0.0 & 1.4 & 1.0 & Mangueta, rótulas, tirante \\
Frenada en curva    & -1.1 & 1.0 & 1.0 & Combinado: suele mandar en la mangueta \\
Aceleración         &  1.0 & 0.0 & 1.0 & Trapecios traseros, palieres \\
Bache con frenada   & -1.0 & 0.0 & 3.0 & \textbf{Casi siempre el crítico} \\
Bordillo vertical   &  0.0 & 0.0 & 3.5 & Pushrod, balancín, amortiguador \\
\bottomrule
\end{tabularx}
\end{center}

\begin{clave}
Los casos combinados no son la suma de los puros, y son los que mandan. Un
neumático no puede dar \SI{1.5}{\gacc} de frenada y \SI{1.4}{\gacc} de lateral a
la vez: la elipse de fricción del apartado \ref{sec:neumatico:03} lo prohíbe. Por
eso el caso de frenada en curva se plantea sobre la elipse ---por ejemplo
\SI{1.1}{\gacc} y \SI{1.0}{\gacc}, que están dentro--- y no sumando los máximos,
que daría cargas imposibles y una suspensión absurdamente pesada.
\end{clave}

% ---------------------------------------------------------------------
\section{Factores de bump y coeficientes de seguridad}
\label{sec:casos_carga:02}
\index{factor de bache}\index{coeficiente de seguridad}

\subsection{El factor de bache}

La carga vertical de un bache no se calcula: se estima con un factor sobre la
estática. Es la parte más arbitraria del capítulo y conviene reconocerlo. Un
factor de \numrange{3}{4} sobre la carga estática es lo habitual en Formula
Student, y lo que hay detrás es que la pista tiene bordillos, juntas y conos,
y que el coche los va a pisar.

Lo que sí se puede hacer es acotarlo por arriba con un argumento físico: la
fuerza vertical máxima que puede transmitir el neumático está limitada por su
rigidez vertical y por el recorrido disponible antes de tocar el tope. Si el
producto de esos dos números da menos que el factor supuesto, el factor sobra.

\begin{ojo}
\textbf{Cuidado con multiplicar factores.} Un caso de \SI{3}{\gacc} vertical con
un coeficiente de seguridad de \num{2} y encima un factor de bache de \num{1.5}
está pidiendo \SI{9}{\gacc} a la pieza. Cada factor tiene que aparecer una sola
vez y estar escrito dónde se aplica. Es el mecanismo por el que las suspensiones
de primer coche acaban pesando el doble de lo necesario.
\end{ojo}

\subsection{Coeficientes de seguridad con criterio}

Un coeficiente de seguridad cubre \emph{incertidumbre}, no ignorancia. Se elige
mirando qué es lo que no se sabe:

\begin{center}
\begin{tabularx}{\textwidth}{@{}lXl@{}}
\toprule
\textbf{Situación} & \textbf{Qué incertidumbre hay} & \textbf{Coef.} \\
\midrule
Carga medida, material certificado, pieza mecanizada & Poca & \numrange{1.3}{1.5} \\
Carga estimada, material conocido & Media & \numrange{1.5}{2.0} \\
Carga estimada, soldadura propia, material sin certificar & Alta & \numrange{2.0}{3.0} \\
Pieza cuyo fallo es catastrófico (dirección, frenos) & La que sea, más margen & \textbf{$\geq$ 2, y ensayar} \\
\bottomrule
\end{tabularx}
\end{center}

\begin{clave}
Un coeficiente de seguridad enorme sobre un caso de carga que nadie ha definido
no protege de nada: solo añade peso y da una sensación falsa de rigor. Es mucho
mejor un coeficiente de \num{1.5} sobre una carga que se sabe de dónde viene que
uno de \num{4} sobre un número inventado. Y ante un juez, la diferencia entre
las dos actitudes se nota en la primera pregunta.
\end{clave}

% ---------------------------------------------------------------------
\section{Cálculo manual previo}
\label{sec:casos_carga:03}
\index{diagrama de sólido libre}

\subsection{Seis barras, seis incógnitas}

Una esquina de doble trapecio con pushrod tiene seis barras biarticuladas: los
cuatro brazos de los trapecios, el tirante de dirección y el pushrod
(\cref{fig:cargas:esquina}). Como son biarticuladas solo trabajan a axial, así
que hay seis incógnitas. Y la mangueta, como sólido rígido en el espacio, da
seis ecuaciones de equilibrio: tres de fuerzas y tres de momentos.

Seis y seis: el problema es \textbf{isostático} y se resuelve exactamente, sin
hipótesis adicionales, resolviendo un sistema lineal $6\times 6$:

\begin{equation}
  \begin{bmatrix} \vect{u}_1 & \cdots & \vect{u}_6 \\
                  \vect{r}_1\times\vect{u}_1 & \cdots &
                  \vect{r}_6\times\vect{u}_6 \end{bmatrix}
  \begin{Bmatrix} N_1 \\ \vdots \\ N_6 \end{Bmatrix}
  = - \begin{Bmatrix} \vect{F} \\ \vect{r}_F\times\vect{F} \end{Bmatrix}
  \label{eq:equilibrio_esquina}
\end{equation}

donde $\vect{u}_i$ es el vector unitario de cada barra, $\vect{r}_i$ la posición
de su extremo en la mangueta y $\vect{F}$ la fuerza en la huella.

\begin{figure}[htbp]
\centering
\begin{tikzpicture}[x=1cm,y=1cm,>={Latex[length=2mm]}]
  \draw[fusalGris] (-0.3,0) -- (7.4,0);
  % rueda
  \draw[fusalGris,fill=fusalGris!12,line width=1pt] (5.9,1.15) circle (1.15);
  \fill[fusalPrim] (5.9,0) circle (2.2pt);
  % mangueta
  \draw[fusalGris,line width=2.5pt] (5.5,0.6) -- (5.35,2.0);
  \draw[fusalGris] (5.42,1.7) -- (4.6,2.85);
  \node[cursonota,anchor=south east] at (4.6,2.85) {mangueta};
  % barras
  \draw[fusalPrim,very thick] (5.5,0.6) -- (1.6,0.9);
  \draw[fusalPrim,very thick] (5.5,0.6) -- (1.9,1.35);
  \draw[fusalPrim,very thick] (5.35,2.0) -- (2.2,1.85);
  \draw[fusalPrim,very thick] (5.35,2.0) -- (2.5,2.25);
  \draw[curvaC,very thick] (5.42,1.25) -- (2.0,1.6);
  \draw[fusalSec,very thick] (5.48,0.85) -- (3.3,3.3);
  \node[cursonota,anchor=south east,text=fusalPrim] at (2.05,0.72)
       {trapecio inferior\\(dos barras)};
  \node[cursonota,anchor=south east,text=fusalPrim] at (2.55,2.3)
       {trapecio superior\\(dos barras)};
  \node[cursonota,anchor=east,text=teal!65!black] at (1.95,1.62) {tirante};
  \node[cursonota,anchor=south,text=fusalSec] at (3.3,3.35) {pushrod};
  % carga en la huella
  \draw[fusalSec,->,very thick] (5.9,0) -- (5.9,1.35);
  \draw[fusalSec,->,very thick] (5.9,0) -- (7.15,0);
  \node[cursonota,anchor=north,text=fusalSec] at (6.6,-0.1) {$\Fy$};
  \node[cursonota,anchor=west,text=fusalSec] at (6.0,1.2) {$\Fz$};
\end{tikzpicture}
\caption[Sólido libre de una esquina]{Diagrama de sólido libre de una esquina.
La mangueta está en equilibrio bajo la fuerza de la huella y las seis fuerzas
axiales de las barras biarticuladas. Seis incógnitas y seis ecuaciones: el
problema es isostático y no hace falta ninguna hipótesis adicional. El pushrod
aparece en el plano por claridad; en el coche va hacia el balancín.}
\label{fig:cargas:esquina}
\end{figure}

\begin{clave}
La comprobación que valida el planteamiento: aplicar \SI{1000}{\newton}
verticales puros y sumar vectorialmente las seis reacciones. Tienen que dar
\SI{1000}{\newton} verticales y cero en las otras dos direcciones, con error de
redondeo. Si no salen, hay un signo o un punto mal. El programa
\texttt{cargas\_esquina.py} hace esa comprobación antes de dar ningún resultado,
y conviene copiar la costumbre en cualquier hoja de cálculo propia.
\end{clave}

\subsection{Lo que sale, y qué manda}

Con la geometría de ejemplo y los seis casos de la tabla, la envolvente por
barra es la de la \cref{fig:cargas:envolvente}. Dos lecturas que conviene
interiorizar:

\begin{itemize}
  \item \textbf{El caso crítico casi siempre lleva bache.} Los casos puramente
        horizontales, que son los que primero se le ocurren a todo el mundo, no
        dimensionan casi nada.
  \item \textbf{El pushrod siempre trabaja a compresión}, y por eso lo limita el
        pandeo y no la resistencia (capítulo \ref{ch:componentes_suspension}). El
        brazo delantero del trapecio inferior, en cambio, ve la mayor tracción,
        porque reacciona el par de frenado como un par entre sus dos brazos.
\end{itemize}

\begin{figure}[htbp]
\centering
\begin{tikzpicture}
\begin{axis}[cursoplot, height=6.0cm, width=0.80\linewidth,
  ybar, bar width=9pt,
  ylabel={carga axial (\si{\newton})},
  symbolic x coords={trap.\ inf.\ del,trap.\ inf.\ tra,trap.\ sup.\ del,
                     trap.\ sup.\ tra,tirante,pushrod},
  xtick=data, x tick label style={rotate=20,anchor=north east,
                                  font=\sffamily\tiny},
  enlarge x limits=0.12,
  ymin=-3200, ymax=6000, ytick={-3000,-1500,0,1500,3000,4500},
  legend style={at={(0.98,0.97)},anchor=north east,
                font=\sffamily\scriptsize,draw=fusalGris!50,fill=white}]
\addplot[fill=fusalPrim!70,draw=fusalPrim] coordinates {
  (trap.\ inf.\ del,5164) (trap.\ inf.\ tra,1327) (trap.\ sup.\ del,413)
  (trap.\ sup.\ tra,1884) (tirante,373) (pushrod,0)};
  \addlegendentry{máxima tracción}
\addplot[fill=fusalSec!60,draw=fusalSec] coordinates {
  (trap.\ inf.\ del,-1060) (trap.\ inf.\ tra,-2383) (trap.\ sup.\ del,-2074)
  (trap.\ sup.\ tra,-695) (tirante,-1402) (pushrod,-2778)};
  \addlegendentry{máxima compresión}
\draw[fusalGris] (axis cs:trap.\ inf.\ del,0) -- (axis cs:pushrod,0);
\end{axis}
\end{tikzpicture}
\caption[Envolvente de cargas por barra]{Envolvente de carga axial de cada barra
sobre los seis casos, calculada con \texttt{cargas\_esquina.py} para la
geometría de ejemplo. Sin coeficiente de seguridad todavía. El pushrod solo
aparece a compresión, que es lo que se espera de un pushrod; si en algún caso
saliera a tracción, habría que revisar la geometría o el caso.}
\label{fig:cargas:envolvente}
\end{figure}

% ---------------------------------------------------------------------
\section{Verificación por FEA}
\label{sec:casos_carga:04}
\index{elementos finitos}\index{convergencia de malla}

El cálculo manual da las cargas de las barras; el elemento finito sirve para las
piezas que \emph{no} son barras: manguetas, balancines, anclajes y soportes,
donde la geometría es tridimensional y no hay fórmula.

\begin{center}
\begin{tabularx}{\textwidth}{@{}lX@{}}
\toprule
\textbf{Para esto sí} & Manguetas, balancines, orejetas y anclajes; concentración
  de tensiones en agujeros y cambios de sección; rigidez de piezas complejas \\
\midrule
\textbf{Para esto no} & Sustituir el cálculo de cargas; obtener el pandeo de un
  tubo, que sale de Euler; convencer a nadie con un dibujo de colores sin
  convergencia demostrada \\
\bottomrule
\end{tabularx}
\end{center}

\begin{ojo}
Un análisis sin \textbf{estudio de convergencia} no vale como verificación. El
procedimiento del capítulo \ref{ch:elementos_finitos} es el mismo aquí: refinar
la malla hasta que la tensión de pico cambie menos de un criterio declarado
---por debajo del \SI{5}{\percent} es razonable--- y \textbf{enseñar la gráfica}.
Sin ella, el número de tensión no significa nada, porque en una singularidad la
tensión crece indefinidamente al refinar.

Y el resultado se contrasta siempre contra algo: la carga de entrada tiene que
coincidir con la del cálculo manual, y las reacciones con las cargas de las
barras. Un FEA que no cuadra con el cálculo a mano tiene un error, y no siempre
está en el cálculo a mano.
\end{ojo}

% ---------------------------------------------------------------------
\section{Ensayos de componente}
\label{sec:casos_carga:05}
\index{ensayo de componente}

No todo se puede ni se debe ensayar. Merece la pena cuando el cálculo tiene
mucha incertidumbre o cuando el fallo es grave:

\begin{center}
\begin{tabularx}{\textwidth}{@{}lXl@{}}
\toprule
\textbf{Ensayo} & \textbf{Qué resuelve} & \textbf{Coste} \\
\midrule
Tracción de un inserto pegado o roscado en tubo & La unión, que es lo que el FEA peor predice & Bajo \\
Compresión de un trapecio completo & Pandeo real con sus extremos reales & Bajo \\
Rigidez de la mangueta & Correlacionar el modelo & Medio \\
Rigidez a torsión del chasis & Cierra el bucle del capítulo \ref{ch:rigidez_torsional} & Medio \\
\bottomrule
\end{tabularx}
\end{center}

\begin{clave}
El ensayo más rentable de todo el bloque es el de \textbf{arrancamiento de un
inserto}. Es donde el cálculo es más flojo ---depende del adhesivo, de la
preparación de la superficie y de la mano que lo hizo---, es donde más fallos
reales aparecen, y se hace con una prensa prestada y media mañana. Si solo se
puede hacer un ensayo, que sea ese.
\end{clave}

% ---------------------------------------------------------------------
%  Cierre del capítulo
% ---------------------------------------------------------------------

\begin{caso}[Manguetas: 37 MPa y lo que eso quiere decir]
Las manguetas se analizaron en Ansys como caso estático estructural bajo cargas
combinadas de peor caso ---frenada fuerte más curva máxima, más las reacciones
del rodamiento y de la dirección---. La tensión de von Mises de pico salió en
torno a \SI{37}{\mega\pascal}, que sobre un aluminio de la serie 6000 deja un
coeficiente de seguridad muy holgado. La convergencia se hizo por refinamiento
h, aceptando menos de un \SI{5}{\percent} de error relativo en la tensión de
pico, y la gráfica de convergencia existe y se enseña.

Hay que saber presentarlo bien, porque un coeficiente de seguridad alto se puede
leer de dos maneras. La nuestra es deliberada: la mangueta se dejó
geométricamente simple, sin vaciados complicados, para que el mecanizado fuera
barato, y el resultado es una pieza más robusta de lo estrictamente necesario.
Es una decisión de coste y de fabricación, no un sobredimensionado accidental, y
así hay que contarlo. Lo que no vale es presentarlo como si hubiéramos
optimizado.

\medskip
\otrodia[Faltan la tabla de casos de carga con las aceleraciones justificadas,
la envolvente de cargas de nuestras barras y la aleación confirmada de la
mangueta.]
\end{caso}

\begin{ojo}
\begin{itemize}
  \item Sumar los máximos de frenada y de curva. La elipse de fricción lo
        prohíbe.
  \item Multiplicar factor de bache por coeficiente de seguridad por factor
        dinámico hasta llegar a cargas que el neumático no puede transmitir.
  \item Dar un coeficiente de seguridad sin decir sobre qué carga está aplicado.
  \item Presentar un FEA sin estudio de convergencia.
  \item No comprobar que las reacciones del FEA coinciden con el cálculo manual.
  \item Dimensionar el pushrod a resistencia cuando lo que lo limita es el
        pandeo.
  \item Olvidar los casos que no son de conducción: montaje, transporte, empujar
        el coche por la carrocería, el gato.
  \item No revisar las cargas cuando cambian las presiones, el neumático o el
        reparto. Los casos de carga caducan.
\end{itemize}
\end{ojo}

\begin{entregable}
\begin{enumerate}
  \item La \textbf{tabla de casos de carga}, con el origen de cada aceleración
        escrito al lado.
  \item La \textbf{hoja de cargas por barra}, con la envolvente de los seis
        casos y el caso crítico de cada una.
  \item El \textbf{informe de FEA} de las piezas no triviales, con estudio de
        convergencia y comprobación contra el cálculo manual.
  \item El \textbf{plan de ensayos} de componente: qué, con qué montaje, con qué
        criterio de aceptación.
\end{enumerate}
\end{entregable}

\begin{juez}
\begin{enumerate}
  \item ¿De dónde salen vuestras cargas?
  \item ¿Por qué esas aceleraciones y no otras?
  \item ¿Cuál es vuestro caso crítico y para qué pieza?
  \item ¿Cómo habéis combinado frenada y curva?
  \item ¿Qué coeficiente de seguridad usáis y sobre qué?
  \item Enseñadme la convergencia de malla.
  \item ¿Habéis ensayado algo físicamente?
  \item ¿Qué pasa con estas cargas si cambiáis de neumático?
\end{enumerate}
\end{juez}
